\documentclass[10pt,a4paper]{amsart}
\usepackage[margin=19mm]{geometry}
\usepackage{amsmath,amssymb,mathtools}
\usepackage[scr=rsfs,cal=euler]{mathalfa}
\usepackage{xfrac}
\usepackage[unicode,hidelinks,bookmarks=false]{hyperref}
\newcommand{\ie}{i.e.\ }
\newcommand{\cf}{cf.\ }
\newcommand{\resp}{respectively\ }
\newcommand{\Subsection}{\subsection}
\providecommand{\nobreakdash}{\nobreak\hbox{-}}
\setlength{\emergencystretch}{2em}
\multlinegap0pt
\usepackage{bm}
\usepackage{bbm}
\usepackage{dsfont}
\usepackage{xspace}
\usepackage{cancel}
\usepackage{mathtools}
\usepackage{amsthm}
\makeatletter
\@ifundefined{lemma}{\newtheorem{lemma}{Lemma}}{}
\@ifundefined{propo}{\newtheorem{propo}{Proposition}}{}
\makeatother
\makeatletter
\expandafter\ifx\csname pr\endcsname\relax
\newenvironment{pr}{{\bf{Proof.~}} \normalfont}
\fi
\makeatother
\title[Spatial SIR epidemic model with varying infectivity in an un]{Spatial SIR epidemic model with varying infectivity in an unbounded domain: Law of Large Numbers}
\author{Armand Kanga, Etienne Pardoux}
\date{Source: arXiv:2511.12350v1 (CC BY 4.0)}
\begin{document}
\maketitle

\noindent\emph{Source and licence.} Spatial SIR epidemic model with varying infectivity in an unbounded domain: Law of Large Numbers. Version arXiv:2511.12350v1 (CC BY 4.0), distributed under the Creative Commons Attribution 4.0 International licence, as recorded in the arXiv licence metadata for this exact version and shown on the abstract page https://arxiv.org/abs/2511.12350v1. Full rights notice: licenses/arxiv-2511.12350-CC-BY-4.0.txt.\par\smallskip

\noindent\textit{Editorial scope.} Counted unit: the Propo whose source label is \texttt{oo} in the article (source0.tex lines 531--546, source label \texttt{oo}), delivered together with the article's own complete original proof of it (source0.tex lines 615--669, one \texttt{proof} environment of 55 lines). No mathematics is written, repaired or re-derived: statement and proof are the article's own text line for line. The proof is detached from the statement in the source: the article states the result at lines 531--546 and prints its proof at lines 615--669, and the proof environment's own header names the statement, so the association is the article's own. Required context retained uncounted: sy (lines 517--528); lem1 (lines 548--551); lem3 (lines 586--593). Every definition, display and label of those lines is retained unchanged. Omitted from this package and not redistributed: 1--516, 529--530, 547--547, 552--585, 594--614, 670--1247. Nothing inside a retained excerpt is omitted or altered: every retained body is delivered exactly as the article's lines are written, so no omission receipt of any kind accompanies this package. Citations: the retained excerpts carry no citation command at all, so no bibliography entry of the article is transcribed and no reference is redistributed. Reference adaptation: none was needed and none was made; every reference inside the delivered unit resolves inside it, so no unresolved reference or citation is produced. Printed slips kept literal: no printed word, symbol, hypothesis or number of the counted statement or of its proof was altered. Layout and class: the article's own class is \texttt{amsart} with packages including latexsym, amssymb, epsfig, bm, amsmath, amsfonts, amsthm, mathrsfs, rotating, appendix, color, multirow, relsize, microtype; the wrapper is the lane's pinned \texttt{amsart} editorial wrapper at 10pt on A4, which restates the theorem-like environments the retained text needs and adds the article's own macro definitions that the retained text needs (none), with extra wrapper packages (bm, bbm, dsfont, xspace, cancel, mathtools). The delivered numbering is the unit's own. Assets: no retained span contains a figure environment, a graphics-inclusion command, a PDF-page inclusion, a table or a TikZ picture; no image file is copied into this package, the package declares no asset dependency and no image file is redistributed with it. Licence: version arXiv:2511.12350v1 of this article is distributed under the Creative Commons Attribution 4.0 International licence, as recorded in the arXiv licence metadata for this exact version and shown on the abstract page https://arxiv.org/abs/2511.12350v1. A full rights notice is delivered at licenses/arxiv-2511.12350-CC-BY-4.0.txt. Characters: every retained excerpt is pure ASCII, so the delivered engine, XeLaTeX with its default Latin Modern text font, reports no missing character.
	\begin{equation}
	\label{sy}
	\left\lbrace
	\begin{aligned}
	(\overline{\mu}_t^{S},\varphi)&= (\overline{\mu}_0^{S},\varphi)- \int_0^t\int_{D}\varphi(x)\overline{\Gamma}(s,x)\overline{\mu}_s^{S}(dx)ds,\\
	(\overline{\mu}_t^{\mathfrak{F}},\varphi) &=\overline{\lambda}^0(t)(\overline{\mu}^I_0,\varphi)+\int_0^t\overline{\lambda}(t-s)\int_{D}\varphi(x)\overline{\Gamma}(s,x)\overline{\mu}_s^{S}(dx)ds,\\
	(\overline{\mu}_t^{I},\varphi)&= (\overline{\mu}_0^{I},\varphi)F_0^c(t)+\int_0^tF^c(t-s)\int_{D} \varphi (x)\overline{\Gamma}(s,x)\overline{\mu}_s^S(dx)ds,\\
	(\overline{\mu}_t^{R},\varphi)&=(\overline{\mu}_0^{R},\varphi)+ (\overline{\mu}_0^{I},\varphi)F_0(t)+\int_0^t F(t-s) \int_{D} \varphi (x)\overline{\Gamma}(s,x)\overline{\mu}_s^S(dx)ds,\\
	\overline{\Gamma}(t,x)&=\int_{D} \dfrac{K(x,y)}{\left[\displaystyle\int_{D}K(z,y)\overline{\mu}(dz) \right]^{\gamma}}\overline{\mu}_t^{\mathfrak{F}}(dy).
	\end{aligned}
	\right.
	\end{equation}	

\begin{lemma}
	\label{lem1}
For any $\varphi \in L^{\infty} (\mathbb{R}^d)$,  $\left \|\displaystyle \int_{D} \Lambda(.,y)\varphi(y)dy\right \|_{\infty}\leq C\|\varphi\|_{\infty}.$
\end{lemma}

\begin{lemma}
	\label{lem3}
	Let $T>0$, and let $(S,\mathfrak{F} )$  be a solution of the first two equations of \eqref{syste}. Then there exists a positive constant $C$ such that:
	\begin{itemize}
		\item $\forall t\in[0;T],$ $\|S(t)\|_{\infty}\leq C;$
		\item $\forall t\in[0;T]$; $\|\mathfrak{F}(t)\|_{\infty}\leq C.$
	\end{itemize}
\end{lemma}

\begin{propo}
	\label{oo}
	The  system \eqref{sy} admits at most one  solution  $\left(\overline{\mu}^S_t,\overline{\mu}_t^{\mathfrak{F}},\overline{\mu}_t^{I},\overline{\mu}_t^{R}, t \geq 0\right)$   which is absolutely continuous with respect to the measure $\overline{\mu}$, with the densities $\left(S(t,.),\mathfrak{F}(t,.),I(t,.) ,R(t,.), t\geq 0\right)$ satisfying for all $x\in D$
	\begin{equation}
	\label{syste}
	\left\lbrace
	\begin{aligned}
	S(t,x)&= S(0,x)- \int_0^t\int_{D} \Lambda(x,y)\mathfrak{F}(s,y)S(s,x)dyds\\
	\mathfrak{F}(t,x) &=\overline{\lambda}^0(t)I(0,x)+\int_0^t\overline{\lambda}(t-s)\int_{D} \Lambda(x,y)\mathfrak{F}(s,y)S(s,x)dyds.\\
	I(t,x)&= I(0,x)F_0^c(t)+\int_0^t F^c(t-s)\int_{D} \Lambda(x,y)\mathfrak{F}(s,y)S(s,x)ds\\
	R(t,x)&= R(0,x)+I(0,x)F_0(t)+\int_0^tF(t-s) \int_{D} \Lambda(x,y)\mathfrak{F}(s,y)S(s,x)dxdyds\\
\Lambda(x,y)&=\dfrac{K(x,y)\overline{\mu}(y)}{\left[\displaystyle\int_{D}K(z,y)\overline{\mu}(dz) \right]^{\gamma}}.
	\end{aligned}
	\right.
	\end{equation}
\end{propo}

\begin{proof}\textit{of Proposition \ref{oo}}.
Step 1: We first show that for all $t\geq 0$ any solution of \eqref{sy}, $\left(\overline{\mu}^S_t,\overline{\mu}_t^{\mathfrak{F}},\overline{\mu}_t^{I},\overline{\mu}_t^{R}\right)$ is absolutely continuous with respect to the measure $\overline{\mu}$, and the densities\\ $\left(S(t,.), \mathfrak{F}(t,.), I(t,.) ,R(t,.)\right)$ verify \eqref{syste}.
 From the first equation of \eqref{sy}, $\overline{\mu}^S_t\leq \overline{\mu}^S_0$. Since $\overline{\mu}^S_0$ is absolutely continuous,  $\overline{\mu}^S_t$ has the same property, and we denote its density by $\overline{\mu}^S(t,x)$. By setting $\overline{\mu}^S(t,x)= S(t,x)\overline{\mu}(x)$, thus $S(t,)$ is the density of $\overline{\mu}^S_t$ with respect to the measure $\overline{\mu}$.\\ 
From the third equation of \eqref{sy}, $\displaystyle \overline{\mu}_t^{I}\leq  \overline{\mu}_0^{I}+ \int_0^t \overline{\Gamma}(s,.)\overline{\mu}^{S}_sds$, thus $\overline{\mu}_t^{I}$ is absolutely continuous , since $\overline{\mu}^I_0$ is absolutely continuous, as well as $\overline{\mu}^{S}_s$ for all s. 
A similar argument applies  to $ \overline{\mu}_t^{\mathfrak{F}}$ and $ \overline{\mu}_t^{R}$.The system of equation \eqref{syste} now follows readily from \eqref{sy}.\\
\smallskip

Step 2: We will verify that  equation \eqref{syste} has at most one solution, which implies that \eqref{sy} has at most one solution.  For that sake, it suffices to show that the system made of the first two equations of \eqref{syste}  has most one solution. The first two equations of the system \eqref{syste} constitute the following system
	\begin{equation}
	\label{sys}
	\left\lbrace
	\begin{aligned}
		S(t,x)&= S(0,x)- \int_0^t\int_{D} \Lambda(x,y)\mathfrak{F}(s,y)S(s,x)dyds,\\
	\mathfrak{F}(t,x) &=\overline{\lambda}^0(t)I(0,x)+\int_0^t\overline{\lambda}(t-s)\int_{D} \Lambda(x,y)\mathfrak{F}(s,y)S(s,x)dyds,\\
 \Lambda(x,y)&=\dfrac{K(x,y)\overline{\mu}(y)}{\left[\displaystyle\int_{D}K(z,y)\overline{\mu}(dz) \right]^{\gamma}}\,.
	\end{aligned}
	\right.
	\end{equation}
	Let $\left(f_1(t,.), g_1(t,.)\right)$ and $\left(f_2(t,.), g_2(t,.)\right)$ be two solutions of the above system with the same initial condition.\\
	On the one hand	
	\begin{align}
	f_1(t,x)-f_2(t,x)&= \int_0^t  \left(f_2(s,x)-f_1(s,x)\right)\int_{D}\Lambda(x,y)g_2(s,y)dyds\nonumber\\
	&+\int_0^tf_1(s,x)\int_{D}\Lambda(x,y)\left( g_2(s,y)-g_1(s,y)\right)dyds\nonumber
	\end{align}
Using the result of Lemma \ref{lem1}, we obtain
	\begin{align}
	\|f_1(t,.)-f_2(t,.)\|_{\infty}&\leq C\int_0^t \|f_2(s,.)-f_1(s,.)\|_{\infty}\|g_2(s,.)\|_{\infty}ds\nonumber\\
	&+ C\int_0^t\| g_2(s,.)-g_1(s,.)\|_{\infty} \|f_1(s,.)\|_{\infty}ds\nonumber
	\end{align}
This combined with Lemma \ref{lem3}, implies that
	\begin{align}
	\label{equa5}
	\|f_1(t,.)-f_2(t,.)\|_{\infty}&\leq C\int_0^t \left(\|f_2(s,.)-f_1(s,.)\|_{\infty}+\| g_2(s,.)-g_1(s,.)\|_{\infty}\right) ds
	\end{align}
	Moreover,
	\begin{align}
	g_1(t,x)-g_2(t,x)&=\int_0^t\overline{\lambda}(t-s) (f_1(s,x)-f_2(s,x))\int_{D}\Lambda(x,y)g_1(s,y)dyds\nonumber\\
	&+\int_0^t \overline{\lambda}(t-s)f_1(s,x)\int_{D}\Lambda(x,y) (g_1(s,y)-g_2(s,y))dyds\nonumber\\
	|g_1(t,x)-g_2(t,x)|&\leq C\int_0^t \| f_1(s,.)-f_2(s,.)\|_{\infty}\left\|\int_{D}\Lambda(x,y)g_1(s,y)dy\right\|_{\infty}ds\nonumber\\
	&+C \int_0^t \left\|\int_{D}\Lambda(x,y) (g_1(s,y)-g_2(s,y))dy\right\|_{\infty}ds\nonumber
	\end{align}
	Using the results of Lemma \ref{lem1} and \ref{lem3}, we obtain
	\begin{align}
	\label{equa6}
\|g_1(t,.)-g_2(t,.)\|_{\infty}	&\leq C\int_0^t \left(\| f_1(s,.)-f_2(s,.)\|_{\infty}+\|g_1(s,.)-g_2(s,.)\|_{\infty}\right)ds
	\end{align}
	From \eqref{equa5} and \eqref{equa6}, we infer that
	\begin{align*}
	\|g_1(t,.)-g_2(t,.)\|_{\infty}+\|f_1(t,.)-f_2(t,.)\|_{\infty}&\leq C \int_0^t  \left(\|g_1(s,.)-g_2(s,.)\|_{\infty}+\|f_1(s,.)-f_2(s,.)\|_{\infty}\right)ds
	\end{align*}
	Using Gronwall's inequality, we conclude that 
	\begin{align*}
	\|g_1(t,.)-g_2(t,.)\|_{\infty}+\|f_1(t,.)-f_2(t,.)\|_{\infty}&=0. 
	\end{align*}
\end{proof}

\end{document}
